\documentclass[12 pt, letterpaper]{hmcpset}
\usepackage{hw-preamble}
\setlist[enumerate, 1]{label = (\roman*)}

\name{Jayden Thadani}
\class{MATH 145 --- Topics in Geometry and Topology}
\assignment{Homework 2}
\duedate{5th February 2025}

\begin{document}

\begin{problem}[10.]
	Recall from class that $f, g \in \operatorname{loops}(D)$ are homotopic, written $f \simeq g$, if there exists a continuous function (called a \emph{homotopy})
	\[
		H : [0, 1] \times [0, 1] \to D
	\]
	such that
	\begin{align*}
		H(0,t) & = f(t) & \text{for all $t \in [0,1]$}, \\
		H(1,t) & = g(t) & \text{for all $t \in [0,1]$}, \\
		H(s,0) & = H(s,1) & \text{for all $s \in [0,1]$}.
	\end{align*}

	Show that $\simeq$ is an equivalence relation.
\end{problem}

\begin{solution}
	Note that a homotopy giving $f \simeq g$ is essentially a continuous function from the cylinder $[0, 1] \times S^{1}$ such that each circular ``section'' of the cylinder is mapped to a loop with $f$ at the bottom and $g$ at the top. We can use this geometric intuition to describe the homotopies that make $\simeq$ an equivalence relation.

	First, the identity homotopy gives $f \simeq f$ with each copy of $S^{1}$ mapped to the loop $f$.

	To be exact, $H : [0, 1] \times [0, 1] \to D$ by $H(s, t) = f(t)$ for all $s \in [0, 1]$. This is continuous since it is constant in $s$ and $f$ is a continuous function of $t$ --- each open set $V$ of $f$ has as its pre-image the open set $[0, 1] \times U$ where $U \subseteq S^{1}$ is the pre-image of $V$ in the map $f : S^{1} \to \mathbb{R}^{2}$. It clearly satisfies $H(0, t) = H(1, t) = f(t)$. It satisfies $H(s, 0) = H(s, 1)$ for each $s$ since $H(s, 0) = f(0) = f(1) = H(s, 1)$.

	Second, given any loops $f \simeq g$ (a cylinder with $f$ at the bottom and $g$ at the top) we can ``flip the cylinder'' to get $g \simeq f$ (a cylinder with $g$ at the bottom and $f$ at the top).

	Formally, if we have a homotopy $H_{1}$ that gives $f \simeq g$, then the homotopy $H_{2}$ with $H_{2}(s, t) = H_{2}(1 - s, t)$ gives $g \simeq f$. It is continuous as the composition of $H_{1}$ with continuous function $(s, t) \mapsto (1 - s, t)$. It also satisfies $H_{2}(0, t) = H_{1}(1 - 0, t) = g(t)$ and $H_{2}(1, t) = H_{1}(1 - 1, t) = f(t)$. Finally, each section is indeed a loop --- $H_{2}(s, 0) = H_{1}(1 - s, 0) = H_{1}(1 - s, 1) = H_{2}(s, 1)$

	Finally given three loops with $f \simeq g$ and $g \simeq h$, we can scale the two cylinders down by half and stack them to get a cylinder with bottom $f$ and top $h$ (with $g$ at the midway point). That is, $f \simeq h$.

	Suppose we are given $f \simeq g$ by a homotopy $H_{1}$ and $g \simeq h$ by another homotopy $H_{2}$. Then consider $H_{3} : [0, 1] \times [0, 1] \to D$ by
	\[
		H_{3}(s, t) = \begin{cases}
			H_{1}(s/2, t) & s \in [0, 1/2] \\
			H_{2}(1/2 + s/2, t) & s \in [1/2, 1].
		\end{cases}
	\]

	It is continuous at each $s \neq 1/2$ since in small enough neighbourhood around $s$, it is one of the two continuous functions $H_{1}$ and $H_{2}$. At $s = 1/2$, $H_{3}$ is continuous since every open neighbourhood containing $H(1/2, t)$ is the (open) gluing (along a segment of the curve) of the open pre-images of the neighbourhood under $H_{1}$ and $H_{2}$. It clearly satisfies $H_{3}(0, t) = H_{1}(0, t) = f(t)$ and $H_{3}(1, t) = H_{2}(1, t) = h(t)$.
\end{solution}

\pagebreak

\begin{problem}[11.]
	Let $p(z) = z^{d} + \sum_{k = 0}^{d - 1} a_{k} z^{k}$ be a monic complex polynomial. Let $R$ be such that
	\[
		R > \lvert a_0 \rvert + \dots + \lvert a_{d-1} \rvert \quad \text{and} \quad R \ge 1.
	\]

	Show that
	\[
		\lvert z^{d} \rvert > \left \lvert \sum_{k = 0}^{d - 1} a_{k} z^{k} \right \rvert
	\]
	for any complex number of magnitude $\lvert z \rvert = R$.
\end{problem}

\begin{solution}
	This follows as a chain of inequalities ---
	\begin{align*}
		\left \lvert \sum_{k = 0}^{d - 1} a_{k} z^{k} \right \rvert & \le \sum_{k = 0}^{d - 1} \lvert a_{k} \rvert \lvert z \rvert^{k} \\
									& \le \sum_{k = 0}^{d - 1} \lvert a_{k} \rvert R^{k} \\
									& \le \sum_{k = 0}^{d - 1} \lvert a_{k} \rvert R^{d - 1} \\
									& \le \left ( \sum_{k = 0}^{d - 1} \lvert a_{k} \rvert \right )  R^{d - 1} \\
									& < R R^{d - 1} \\
									& = R^{d}.
	\end{align*}
	Here, the first inequality is the triangle inequality, and the third holds because $R \ge 1$.
\end{solution}

\pagebreak

\begin{problem}[12.]
	\begin{enumerate}
		\item Let $f, g \in \operatorname{loops}(\mathbb{C})$ such that $\lvert f(t) \rvert > \lvert g(t) \rvert$ for all $t \in [0, 1]$. Show that $f$ and $f + g$ both belong to $\operatorname{loops}(\mathbb{C} \setminus \{ 0 \})$, and that $f \simeq f + g$ in $\mathbb{C} \setminus \{ 0 \}$.
		\item Give an example of $f, g \in \operatorname{loops}(\mathbb{C})$ such that $f$ and $f + g$ both belong to $\operatorname{loops}(\mathbb{C} \setminus \{ 0 \})$ but $w(f, 0) \neq w(f + g, 0)$. Identify a value of $t$ for which $\lvert f(t) \rvert \le \lvert g(t) \rvert$.
	\end{enumerate}
\end{problem}

\begin{solution}
	\begin{enumerate}
		\item Since $\lvert f(t) \rvert > \lvert g(t) \rvert \ge 0$ for all $t$ $\lvert f(t) \rvert > 0$ and thus, $f \in \operatorname{loops}(\mathbb{C} \setminus \{ 0 \})$. Also, by the reverse triangle inequality, $\lvert (f + g)(t) \rvert > \lvert f(t) \rvert - \lvert g(t) \rvert = 0$. That is, $f + g$ is in $\operatorname{loops}(\mathbb{C} \setminus \{ 0 \})$.

			For each $t \in [0, 1]$, the support of the line segment $[f(t), (f + g)(t)]$ is $\{ f(t) + s g(t) \mid s \in [0, 1] \} \subset \mathbb{C}$. In fact (again by the reverse triangle inequality) $\lvert f(t) + s g(t) \rvert \ge \lvert f(t) \rvert - s \lvert g(t) \rvert > 0$. Thus, support of each line segment is in fact a subset of $\mathbb{C} \setminus \{ 0 \}$.

			Thus, by the straight line homotopy, $f \simeq f + g$ in $\mathbb{C} \setminus \{ 0 \}$.
		\item Consider $f = \{ z \in \mathbb{C} \setminus \{ 0 \} \mid \lvert z - 1/2 \rvert = 1/4 \}$, the circle of radius $1/4$ centred at $1/2$, and $g = \{ z \in \mathbb{C} \setminus \{ 0 \} \mid \lvert z \rvert = 1 \}$, the circle of radius $1$ centred at $0$.

			$f$ winds about $0$ zero times but $f + g$ winds about $0$ once. That is, $w(f, 0) \neq w(f + g, 0)$

			$\lvert f(t) \rvert \le 3/4 < \lvert g(t) \rvert = 1$ for all $t$, including, say, $0$.
	\end{enumerate}
\end{solution}

\pagebreak

\begin{problem}[13. (Proof of the fundamental theorem of algebra)]
	Let $p(z) = z^{d} + \sum_{k = 0}^{d - 1} a_{k} z^{k}$ be a complex polynomial. Let $R$ be a real number such that:
	\[
		R > \lvert a_{0} \rvert + \dots + \lvert a_{d - 1} \rvert \quad \text{and} \quad R \ge 1
	\]

	Define five loops in $\mathbb{C}$ as follows:
	\begin{align*}
		f(t) & = e^{2\pi i d t} \\
		g(t) & = R^d e^{2\pi i d t} \\
		h(t) & = p(Re^{2\pi i t}) \\
		k(t) & = a_0 \\
		\ell(t) & = 1
	\end{align*}

	Then $f, g, \ell$ self-evidently are loops in $\mathbb{C} \setminus \{ 0 \}$ and problem 12 implies that $h$ is too. Finally, if we assume that $p$  has no roots then $a_{0} = 0$ and so the same is true for $k$. The next step is to show that these loops are homotopic in $\mathbb{C} \setminus \{ 0 \}$.
	\begin{enumerate}
		\item Show that $f \simeq g$.
		\item Show that $g \simeq h$.
		\item Assume that $p$ has no roots. Show that $h \simeq k$.
		\item Assume that $p$ has no roots. Show that $k \simeq \ell$.
	\end{enumerate}

	These four results lead to a contradiction if we assume that $p$ has no roots, because then the continuous winding theorem implies:
	\[
		d = w(f, 0) = w(g, 0) = w(h, 0) = w(k, 0) = w(\ell, 0) = 0.
	\]
	It follows that $p$ must have at least one root.
\end{problem}

\begin{solution}
	\begin{enumerate}
		\item Notice that $(g - f)(t) = (R^{d} - 1) e^{2 \pi i d t}$. Since $R > 1$ this means $\lvert (f - g)(t) \rvert = \lvert (g - f)(t) \rvert < \lvert f(t) \rvert$ for all $t$. By problem 12 part (i), this means $(f - g) + g = f \simeq g$ (in $\mathbb{C} \setminus \{ 0 \}$).
		\item The statement of problem 11 is exactly that $\lvert g(t) \rvert = R^{d} > \lvert g(t) - h(t) \rvert$. Thus, $\lvert h(t) - g(t) \rvert < \lvert g(t) \rvert$ too. Thus, $(h - g) + g = h \simeq g$.
		\item We can use the ``pushforward homotopy'' --- the fact that continuous functions preserve homotopy. If $p$ has no roots, then it is a continuous function $\mathbb{C} \to \mathbb{C} \setminus \{ 0 \}$. Note that $k = a_{0} = p(0)$. Since the circle of radius $R$ (the loop $R e^{2 \pi i t}$) is homotopic to the constant loop $0$ in $\mathbb{C}$, the loop obtained from the continuous function $p$ applied to the circle (the loop $h$) is homotopic to the constant loop $k = p(0)$ in $\mathbb{C} \setminus \{ 0 \}$.
		\item Since $\mathbb{C} \setminus \{ 0 \}$ is path connected, any two constant loops are homotopic --- given a path $\gamma$ from $a_{0}$ to $1$, there is a homotopy $H(s, t) = \gamma(s)$ giving $k \simeq \ell$.
	\end{enumerate}
\end{solution}

\pagebreak

\begin{problem}[14.]
	Consider the complex polynomial
	\[
		p(z) = 9 + 3z + 22z^2 - 17z^3 + 3z^4.
	\]

	We can deduce from 13 (ii) that
	\[
		w(p(Re^{2\pi i t}), 0) = 4
	\]
	when $R$ is large enough. Thus, $p$ applied to a large circle has a winding number of $4$.

	Now we will consider the winding number of $p$ applied to a small circle around a root.

	\begin{enumerate}
		\item Show that $p(z)$ has a double root at $z = 3$, and determine all four roots of $p$.
		\item What are the roots of the polynomial $q(z) = p(3 + z)$?
		\item Write out $q(z) = p(3+z)$ in the form $a_2 z^2 + a_3 z^3 + a_4 z^4$.
		\item Let $r > 0$ be chosen sufficiently small that
			\[
				(\lvert a_{3} \rvert + \lvert a_{4} \rvert) \, r < |a_2| \quad \text{and} \quad r \leq 1.
			\]
			Show that
			\[
				w(p(3 + re^{2\pi i t}), 0) = 2.
			\]

			Write down a specific such value of $r$.
		\item Conjecture a lemma that generalizes (iv) to an arbitrary polynomial and a small loop around an arbitrary point in the plane. You are not required to prove the lemma.
		\item Draw the four roots of~$p$ in the complex plane, and draw two loops
			\[
				\gamma_{1}, \gamma_{2} \in \operatorname{loops}(\mathbb{C} \setminus \{ \text{roots of } p \})
			\]
			such that
			\[
				w(p \circ \gamma_{1}, 0) = -4 \quad \text{and} \quad w(p \circ \gamma_{2}, 0) = 7.
			\]

	\end{enumerate}
\end{problem}

\begin{solution}
	\begin{enumerate}
		\item We will show that $(z - 3)^{2}$ divides $p$. Specifically, consider
			\begin{align*}
				(z - 3)^{2} (3 z^{2} + z + 1) & = (z^{2} - 6 z + 9) (3z^{2} + z + 1) \\
							      & = 3 z^{4} - 18 z^{3} + z^{3} + z^{2} + 27 z^{2} - 6 z^{2} - 6z + 9z + 9 \\
							      & = 3 z^{4} - 17 z^{3} + 22 z^{2} + 3 z + 9 \\
							      & = p(z).
			\end{align*}

			From the quadratic formula, we see that $3z^{2} + z + 1$ has two roots ---
			\[
				z = \frac{-1 \pm \sqrt{11} i}{6}
			\]
			each of multiplicity one. Together with the double root at $z = 3$, these account for all of the ``four'' roots of $p$.
		\item $z$ is a root of $q$ exactly when $z + 3$ is a root of $p$. Thus, $q$ has a double root at $0$, and simple roots at $(-4 \pm \sqrt{11} i)/6$.
		\item
			\[
				\boxed{
					q(z) = 31 z^{2} + 19 z^{3} + 3 z^{4}.
				}
			\]

			
			\begin{align*}
				q(z) & = ((3 + z) - 3)^{2} (3 (3 + z)^{2} + (3 + z) + 1) \\
				     & = z^{2} (3 z^{2} + 19 z + 31) \\
				     & = 31 z^{2} + 19 z^{3} + 3 z^{4}.
			\end{align*}
		\item Denote by $f$ the loop $t \mapsto p(3 + r e^{2 \pi i t})$ and by $g$ the loop $t \mapsto a_{2} r e^{2 \pi i d t}$ for $d = 2$. Since $r \le 1$, for all $z \in \mathbb{C}$ with $\lvert z \rvert \le r$
			\begin{align*}
				\lvert a_{3} z^{3} + a_{4} z^{4} \rvert & \le \lvert a_{3} \rvert r^{3} + \lvert a_{4} \rvert r^{4} \\
									& < ( \lvert a_{3} \rvert + \lvert a_{4} \rvert ) r^{3} \\
									& \le \lvert a_{2} \rvert r^{2}.
			\end{align*}

			That is, $\lvert g(t) \rvert > \lvert (f - g)(t) \rvert$. By problem 12 part (i), $(f - g) + g = f \cong g$. Since $g$ has (by definition) winding number $2$ about $0$, so does $f$. That is, $w(p(3 + r e^{2 \pi i t}), 0) = 2$. 
		\item
			\begin{lemma}
				Given a polynomial $p \in \mathbb{C}[z]$, there exists a loop $f$ around $z \in \mathbb{C}$ of small enough radius distance from $z$ such that $w(p \circ f, 0)$ is the multiplicity of the polynomial's root at $z$.
			\end{lemma}
		\item
			We can construct the loops $\gamma$ as winding around each of the roots an integer number of times. The total winding number of the loop $p \circ \gamma$ is the sum of the number of times it winds around each root, weighted by the multiplicity of the root.
			\begin{figure}[H]
				\centering
				\includegraphics[width = 0.4 \textwidth]{figures/P14vi gamma1.pdf}
				\includegraphics[width = 0.4 \textwidth]{figures/P14vi gamma2.pdf}
				\caption{$\gamma_{1}$ winds around all the roots once, clockwise. $\gamma_{2}$ winds around $3$ and $(-1 + \sqrt{11} i) / 6$ twice, and winds around $(-1 - \sqrt{11} i) / 6$ once.}
			\end{figure}
	\end{enumerate}
\end{solution}

\end{document}
